The existing multisignature service decides whether to intervene.
It can observe regional health, the current export session, accepted
DATA progress and operational reports. Missing-block and delivery-rate
tests belong to this policy. They are not consensus-layer transitions
or additional certificate predicates.

The interface is:
\[
 \text{verified observations and reports}
 \ \longrightarrow\
 \text{future L1 STATUS and MSIG records}.
\]
Other participants may run monitors and send reports. Only finalized,
authorized L1 controls can affect public eligibility or cause Force.
A report delivered to one region cannot change that region's
independent replay result.

To exclude a failed region and force a handoff, the controller explicitly
schedules OFF and MSIG. To force a handoff without changing eligibility,
it schedules MSIG alone. ON restores future selection eligibility.
There is no mandatory restoration lease, service credit or cross-region
fairness obligation.

The certifier checks the control's authorization, source completeness,
activation position and target version. It does not verify whether an
internal missing-block threshold was met. A committee instruction has
the same external meaning regardless of the report that motivated it.
The committee cannot supply a preferred retained block; replay uses
the actual pre-payload $G$.

This version leaves detection thresholds, report quality, signing policy,
response time and intervention frequency to a separate governance design.
Continuous Hard Force is permitted, even during preparation or before
the new session has produced DATA. Such interference can prevent
progress without changing already accepted output.

If a region is offline and no valid MSIG is submitted, the consensus
layer does not force a switch. The L1 service bound starts at actual
submission; it excludes detection, decision and signature-collection
delay. Recovery liveness therefore needs the separate committee
response and service-opportunity conditions of \cref{sec:progress}.
